% ERL Letter: Introduction (D-107). Every number is a \metric{} macro; every citation key is
% in refs/cite_list.csv; no citation markers remain.

Closed seas and saline lakes are shrinking\citep{wurtsbaugh2017,yao2023}, and many inland seas and lakes will keep falling through this century as the climate warms\citep{prange2020}. Where the floor is shallow and flat, each metre of fall exposes a wide band of seabed, whose position decides which harbours, channels and habitats are lost. Projections of the emerging seabed combine a water-level pathway with the shape of the basin floor, and in recent assessments that shape comes from a global bathymetry grid\citep{court2025,duan2025}. On a nearly flat shelf a small vertical error in the grid moves its predicted shoreline a long way, here about \metric{rev.w.r.km_per_m_min}--\metric{rev.w.r.km_per_m_max}~km per metre of level. The grids are compiled from sparse soundings: only about \metric{lit.weatherall2015.control_share_pct}\% of the non-land cells of the 2014 grid of the General Bathymetric Chart of the Oceans (GEBCO) were based on bathymetric control data, which include pre-prepared grids that may contain interpolated values\citep{weatherall2015}, and in the current GEBCO grid the North Caspian seabed comes from such a pre-generated grid of mixed sources, whose vertical reference is not documented\citep{gebco2026}. Yet the grids are rarely tested against the seabed that has already emerged.

By the end of 2024 the level of the Caspian Sea had fallen to about \metric{lit.ginzburg2025.level_end2024_bhs_m}~m in the Baltic Height System, below the 1977 minimum of \metric{lit.ginzburg2025.min1977_bhs_m}~m, after falling by \metric{lit.ginzburg2025.fall_rate_2020_2024_cm_per_yr}~cm per year over 2020--2024\citep{ginzburg2025}. One study using \metric{lit.hoseini2025.n_gcms} climate models gives median changes from 2021 to 2100 of \metric{lit.hoseini2025.medium_median_m}~m for a medium-emission scenario and \metric{lit.hoseini2025.high_median_m}~m for a high-emission one\citep{hoseini2025}. Another puts the fall over 2020--2100 at about \metric{lit.samant2023.ssp245_m}~m under SSP2-4.5 and \metric{lit.samant2023.ssp585_m}~m under SSP5-8.5, each including about \metric{lit.samant2023.trend_2100_m}~m from the observed 1996--2021 trend carried forward\citep{samant2023}. Much of the seabed that emerges will lie in the North Caspian, a shelf that holds less than \metric{lit.akbari2020.nc_volume_share_pct}\% of the sea's volume\citep{akbari2020}, with a mean depth of \metric{rev.w.v.depth2026_m_min}--\metric{rev.w.v.depth2026_m_max}~m at the 2026 level. Assessments of what the fall will expose use global grids: one, based on GEBCO\_2022, found that a decline of \metric{lit.court2025.decline_lo_m}--\metric{lit.court2025.decline_hi_m}~m would reduce existing marine protected area coverage by up to \metric{lit.court2025.mpa_loss_max_pct}\%\citep{court2025}; another took the sea's topography from GEBCO and concluded that a fall to \metric{lit.duan2025.level_2100_m}~m would remove the northern Caspian\citep{duan2025}. These grids have been checked against the emerged seabed only qualitatively\citep{akbari2020}. A whole-sea comparison found satellite areas within \metric{lit.bocharov2026.max_area_diff_pct}\% of those from historical area--level estimates and a chart-digitised elevation model, none of them a global grid\citep{bocharov2026}; in a whole-sea total, errors of opposite sign on different coasts can cancel.

A falling closed sea surveys its own floor. Each cloud-free satellite image records where the water's edge lay, and satellite altimetry records how high the water stood that day, so every dated waterline is a contour of the seabed\citep{schroder2022,li2022,schwatke2020}; lake hypsometry has been mapped in this way, or from flood frequency and level series, within the observed levels\citep{cretaux2016,busker2019,fassoniandrade2020}, and the bathymetry of shallow lakes from water occurrence and satellite heights\citep{armon2020}, much as waterlines at known tide heights map intertidal topography\citep{sagar2017,bishoptaylor2019}. Since \metric{levels.dahiti.smoothed.date_of_max} the Caspian has fallen by \metric{levels.dahiti.smoothed.range_m}~m, and its waterline has swept across the North Caspian shelf (figure~\ref{fig:erl_thennow}); a satellite shoreline dataset records that retreat since 2016 but assigns no heights to the shorelines\citep{dedova2026a}. We ask whether five global grids, flooded from the open sea at each observed level, reproduce the observed shorelines; how the same audit treats them at Great Salt Lake, where a lidar-based surface can serve as a control; and what they imply for the next metres of fall.
